\begin{helper}
Let $G$ be a finite graph, let $A$ be an activation set, and fix an ordered
triple of vertices $(a,b,c)$.  Put $Q=\{a,b,c\}$.  Partition the vertices of
$A\setminus Q$ according to their nonempty adjacency pattern to the ordered
triple, and write
\[
n_a,n_b,n_c,n_{ab},n_{ac},n_{bc},n_{abc}
\]
for the seven corresponding cardinalities.  Then the exponents which occur in
the ordered chamber with priority order $a,b,c$ satisfy
\[
n_a+n_{ab}+n_{ac}+n_{abc}
=|\{w\in A\setminus Q:w\sim a\}|,
\]
\[
n_b+n_{bc}
=|\{w\in A\setminus Q:w\not\sim a,\ w\sim b\}|,
\]
and
\[
n_c
=|\{w\in A\setminus Q:w\not\sim a,\ w\not\sim b,\ w\sim c\}|.
\]
\end{helper}
\begin{proof}
The statement is a disjoint finite-set bookkeeping identity.  For the first
identity, consider the set
\[
S_a=\{w\in A\setminus Q:w\sim a\}.
\]
Splitting $S_a$ according to whether $w\sim b$ and then, inside each part,
according to whether $w\sim c$, gives four disjoint classes:
\[
\{a\text{-only}\},\quad \{a,b\text{ only}\},\quad
\{a,c\text{ only}\},\quad \{a,b,c\}.
\]
Their cardinalities are respectively
$n_a,n_{ab},n_{ac},n_{abc}$, so
\[
|S_a|=n_a+n_{ab}+n_{ac}+n_{abc}.
\]

For the second identity, consider
\[
S_b=\{w\in A\setminus Q:w\not\sim a,\ w\sim b\}.
\]
Splitting $S_b$ according to whether $w\sim c$ gives the two disjoint classes
$\{b\text{-only}\}$ and $\{b,c\text{ only}\}$, with cardinalities
$n_b$ and $n_{bc}$.  Hence $|S_b|=n_b+n_{bc}$.

The third identity is exactly the definition of the $c$-only class:
\[
\{w\in A\setminus Q:w\not\sim a,\ w\not\sim b,\ w\sim c\}.
\]
Thus its cardinality is $n_c$.
\end{proof}
